\documentclass[11pt,a4paper]{article}
\usepackage[margin=24mm]{geometry}
\usepackage{amsmath,amssymb}
\input{glyphtounicode}
\pdfgentounicode=1
\usepackage[hidelinks]{hyperref}
\setlength{\parindent}{0pt}
\setlength{\parskip}{7pt}
\newcommand{\der}[1]{\frac{\mathrm d #1}{\mathrm d\alpha}}
\begin{document}
\begin{center}
{\Large How the two lines of equations 15 and 16 connect}\\[5pt]
{\small Clarification of the urban compact methods report}
\end{center}

The individual flows enter through a \textbf{total-allocation constraint}:
lane additions must sum to the admitted mass, and leg discharge rates must
sum to the usable output rate. Substituting the lower cases into the
derivative of that constraint produces the upper equation.

\textbf{Equation 15: lane additions share one mass budget.}
Here $x_\ell$ is the mass of the current packet actually admitted to lane
$\ell$. It determines how much traffic enters that lane; it is not a
discharge rate. The total admitted mass is $V$.
The starting load is $h_\ell$, the admission bound is $u_\ell$, and
$\lambda$ is the common filling level.

On an unchanged set of allocation categories, the forward rule is
\[
 x_\ell=
 \begin{cases}
 \lambda-h_\ell,&\ell\in\mathcal A\quad(0<x_\ell<u_\ell),\\
 u_\ell,&\ell\in\mathcal U\quad\text{(upper bound binds)},\\
 0,&\text{lane remains without an addition},
 \end{cases}
 \qquad \sum_\ell x_\ell=V.
\]
Thus $x_\ell$ matters even though it does not appear explicitly in the
upper gradient formula: its forward value identifies which case applies,
and its derivative must respect the total mass budget.

Differentiate the total, then replace each individual derivative by the
appropriate lower case of equation 15:
\[
 \underbrace{\der{V}=\sum_\ell\der{x_\ell}}_{\text{differentiated total-mass constraint}}
 =\sum_{\ell\in\mathcal A}
       \left(\der{\lambda}-\der{h_\ell}\right)
  +\sum_{\ell\in\mathcal U}\der{u_\ell}.
\]
Lanes remaining at zero contribute zero. Each lane in $\mathcal A$
contains the same unknown $\mathrm d\lambda/\mathrm d\alpha$, so collecting
those terms gives
\[
 |\mathcal A|\der{\lambda}
 =\der{V}-\sum_{\ell\in\mathcal U}\der{u_\ell}
             +\sum_{\ell\in\mathcal A}\der{h_\ell}.
\]
Dividing by $|\mathcal A|$ gives the upper line of equation 15.
\textbf{You substitute the lower cases into the total constraint, rather
than substituting the forward values of $x_\ell$ into the upper line.}

For calculation, first solve the forward allocation and identify its
categories. Given the input sensitivities of $V$, $u_\ell$ and $h_\ell$,
use the upper line to obtain the filling-level sensitivity. Then use the
lower line to obtain each lane's admitted-mass sensitivity. Their sum
equals the sensitivity of $V$ by construction.

\newpage
\textbf{Equation 16: inbound legs share one output-rate budget.}
Here $r_{ij}$ is the discharge rate granted to inbound leg $i$ toward
output $j$. The lane-sharing calculation divides this grant among that
leg's ready lanes; their rates determine transferred mass over each
head-service interval. The total grant toward $j$ is $W_j$.

For a fixed output $j$, the cap of leg $i$ is $C_{ij}$, its green weight
is $w_{ij}$, and $\theta_j$ is the common multiplier of those weights.
With $\mathcal C$ the capped legs and $\mathcal F$ the uncapped legs,
\[
 r_{ij}=
 \begin{cases}
 C_{ij},&i\in\mathcal C,\\
 \theta_jw_{ij},&i\in\mathcal F,
 \end{cases}
 \qquad \sum_i r_{ij}=W_j.
\]
Again, differentiate the total constraint and substitute the lower cases
of equation 16. The uncapped case uses the product rule:
\begin{align*}
 \der{W_j}
 &=\sum_i\der{r_{ij}}\\
 &=\sum_{i\in\mathcal C}\der{C_{ij}}
   +\sum_{i\in\mathcal F}
       \left(w_{ij}\der{\theta_j}
             +\theta_j\der{w_{ij}}\right).
\end{align*}
Collect the terms containing the same multiplier sensitivity:
\[
 \left(\sum_{i\in\mathcal F}w_{ij}\right)\der{\theta_j}
 =\der{W_j}-\sum_{i\in\mathcal C}\der{C_{ij}}
          -\theta_j\sum_{i\in\mathcal F}\der{w_{ij}}.
\]
Dividing by the total uncapped weight gives the upper line of equation 16.
The individual rates disappear explicitly because their lower-case
expressions have replaced them inside their sum.

The calculation again runs from the common quantity to the individual
ones: solve the forward rates and identify the capped legs; use input
sensitivities of total service, caps and weights to calculate the upper
line; then calculate each leg's rate sensitivity using the lower line.

\textbf{Why the two lines are not circular.}
In the derivation, individual responses are inserted into a shared budget
to solve for one unknown common-level response. In numerical evaluation,
that common-level response is calculated first and then distributed back
to the individual flows. Neither upper line requires the individual flow
derivatives as an additional unknown input.

These formulas apply while the categories remain unchanged and there is
at least one uncapped contributor (with positive total weight in equation
16). If none remains, capped flows follow their caps directly; there is
no need to determine a common-level derivative by division.
\end{document}
